\section{Procedimentos metodológicos}
\label{sec:metodologia}

A pesquisa utiliza a PNAD Contínua para construir um painel de transições entre entrevistas consecutivas. A unidade de análise é o par pessoa-trimestre, formado mediante regras de identificação e consistência dos registros. Os procedimentos incluem a preparação dos insumos, a correspondência entre classificações ocupacionais e o cálculo dos indicadores de transição. A documentação dessas etapas busca permitir a reprodução da análise e a identificação de suas limitações.

O desenho relaciona a exposição ocupacional à inteligência artificial aos desfechos de mudança de ocupação, saída do emprego e passagem para a informalidade. As estimações consideram pesos amostrais, controles e efeitos fixos, conforme a especificação documentada no projeto. O marco temporal organiza a comparação entre períodos, mas não representa uma medida de adoção efetiva da tecnologia. A interpretação deverá distinguir associações condicionais de efeitos causais e discutir possíveis mudanças na coleta e na composição da amostra.

Este capítulo é um esqueleto provisório para o exercício ia-6.1. A versão final deverá apresentar as regras de inclusão, as variáveis, a especificação econométrica e as verificações realizadas, sem omitir limitações ou etapas necessárias à reprodução.